%% Illustration de Parseval-Plancherel sur le signal porte : |x(t)|^2 =
%% Pi_theta(t) a gauche, |X(f)|^2 = theta^2 sinc^2(pi f theta) a
%% droite, meme aire (energie). Repris du poly (section 2). Consomme
%% par slides.
\begin{tikzpicture}[x=1.2cm,y=1.2cm]
    \draw[-Latex] (-2.2,0) -- (2.2,0) node[right] {$t$};
    \draw[-Latex] (0,-0.2) -- (0,1.5) node[above] {$|x(t)|^2$};
    \fill[niceblue,fill opacity=0.3] (-0.5,0) rectangle (0.5,1);
    \draw[niceblue,thick] (-2,0) -- (-0.5,0) -- (-0.5,1) -- (0.5,1) -- (0.5,0) -- (2,0);
    \node[below] at (-0.5,0) {$-\theta/2$};
    \node[below] at (0.5,0) {$\theta/2$};
    \node[left] at (0,1) {$1$};
\end{tikzpicture}
\hspace{1cm}
\begin{tikzpicture}
    \begin{axis}[
        width=6.5cm, height=4cm,
        xlabel={$f$}, ylabel={$|X(f)|^2$},
        xmin=-4, xmax=4, ymin=0,
        axis lines=middle,
        xlabel style={at=(current axis.right of origin), anchor=west},
        ylabel style={at=(current axis.above origin), anchor=south},
    ]
    \addplot[domain=-4:4, samples=200, niceblue, thick, fill=niceblue, fill opacity=0.3] {(sin(deg(3.14159*x))/(3.14159*x))^2} \closedcycle;
    \end{axis}
\end{tikzpicture}
